% Generated from locale/tr/content/proof-theory/sequent-calculus/rules-proofs.tex; do not hand-edit.


\olfileid[tr]{pt}{seq}{rp}

\olsection{Kurallar ve \usetoken{P}{proof}}

Bazı tanımlar verip terminolojiyi sabitleyerek başlayalım.

Örneğin \Log{G3c}'de, $!A$ ve~$!B$ içeren sekantlardan $!A \land !B$
içeren sekantları !!{proving}{}ya olanak veren şu iki kuralı ele alalım:
\[
\Axiom$ !A, !B, \Gamma \fCenter \Delta$
\RightLabel{\LeftR{\land}}
\UnaryInf$ !A \land !B, \Gamma \fCenter \Delta$
\DisplayProof
\qquad
\Axiom$\Gamma \fCenter \Delta, !A$
\Axiom$ \Gamma \fCenter \Delta, !B$
\RightLabel{\RightR{\land}}
\BinaryInf$ \Gamma \fCenter \Delta, !A \land !B $
\DisplayProof
\]
Çizginin üstündeki sekantlara \emph{öncüller}, altındaki sekanta
\emph{sonuç} denir. Her kuraldaki $\Gamma$ ve $\Delta$ içinde bulunan
!!{formula}s \emph{bağlam} ya da \emph{yan !!{formula}s} diye adlandırılır.
Sonuçta yer alıp bağlamın parçası olmayan !!{formula}, \emph{başlıca}
!!{formula}; öncüllerde yer alıp bağlamda bulunmayan !!{formula}s ise
\emph{etkin} (ya da \emph{yardımcı}) !!{formula}s diye adlandırılır.

Niceleyici kuralları biraz daha karmaşıktır. Örneğin \Log{G1c}'de
$\lforall$ kuralları şöyledir:
\[\Axiom$ !A(t), \Gamma \fCenter \Delta$
\RightLabel{\LeftR{\lforall}}
\UnaryInf$ \lforall[x][!A(x)],\Gamma \fCenter \Delta$
\DisplayProof
\qquad
\Axiom$ \Gamma \fCenter \Delta, !A(a) $
\RightLabel{\RightR{\lforall}}
\UnaryInf$ \Gamma \fCenter \Delta, \lforall[x][!A(x)]$
\DisplayProof
\]
\LeftR{\lforall} kuralında etkin !!{formula}, $t$ kapalı bir terim, yani
değişkensiz bir terim olmak üzere~$!A(t)$'dir. Böyle bir çıkarımın doğru,
yani şematik \LeftR{\lforall} kuralına uygun olması için; sonuçtaki
başlıca !!{formula} $\lforall[x][!A]$ ve öncüldeki etkin !!{formula}
$\Subst{!A}{t}{x}$ olacak biçimde bir !!{formula}~$!A$, bir değişken~$x$
ve bir kapalı terim~$t$ bulunmalıdır. \RightR{\lforall} kuralı da aynı
biçimde işler; ancak gelişigüzel bir $t$ terimi yerine öncüldeki etkin
!!{formula}, $a$ bir !!a{constant} olmak üzere $\Subst{!A}{a}{x}$
olmalıdır. Bu sabite \RightR{\lforall} çıkarımının \emph{özdeğişkeni}
denir. Çıkarım ancak alt sekant $a$'yı içermiyorsa, yani $a$; $\Gamma$,
$\Delta$ ya da $!A$ içinde geçmiyorsa doğrudur. Buna
\emph{özdeğişken koşulu} denir.

\Log{G1c} ve \Log{G3c} gibi sekant sistemleri, şematik olarak belirtilmiş
bu tür kuralların bir derlemiyle, aksiyom sayılan (yine şematik olarak
belirtilmiş) bazı sekantlardan oluşur. Böyle bir sistemdeki !!^{proof}s,
aksiyomlardan çıkarım kuralları kullanılarak tümevarımsal biçimde üretilen
sonlu sekant ağaçlarıdır. Her yaprak bir aksiyomdur; diğer her sekant ise
hemen üstündeki sekantlardan sistemin çıkarım kurallarından biriyle çıkar.

\begin{defn}\ollabel{defn:proof}
Bir sekant sistemindeki \emph{!!^{proof}s}, tümevarımsal olarak aşağıdaki
gibi tanımlanan sonlu sekant ağaçlarıdır:
\begin{enumerate}
  \item\ollabel{defn:prf:axiom} Her aksiyom bir !!a{proof}{}tır.
  \item\ollabel{defn:prf:one-prem} $\pi_1$, $S_1$ sekantıyla biten bir
  !!a{proof} ve $S$, tek öncüllü bir~$R$ kuralıyla $S_1$'den çıkan bir
  sekantsa
  \[
  \AxiomC{}
  \RightLabel{$\pi_1$}
  \DeduceC{$S_1$}
  \RightLabel{$R$}
  \UnaryInfC{$S$}
  \DisplayProof
  \]
  bir !!a{proof}{}tır.
  \item\ollabel{defn:prf:two-prem} $\pi_1$ ile $\pi_2$, sırasıyla $S_1$
  ve~$S_2$ sekantlarıyla biten !!{proof}s ve $S$, iki öncüllü bir~$R$
  kuralıyla $S_1$ ile $S_2$'den çıkan bir sekantsa
  \[
  \AxiomC{}
  \RightLabel{$\pi_1$}
  \DeduceC{$S_1$}
  \AxiomC{}
  \RightLabel{$\pi_2$}
  \DeduceC{$S_2$}
  \RightLabel{$R$}
  \BinaryInfC{$S$}
  \DisplayProof
  \]
  bir !!a{proof}{}tır.
\end{enumerate}
\end{defn}

\Olref{defn:proof}, !!{proof}s{}ı sekant ağaçları olarak tanımlar. Bu
ağaçların ``yukarı doğru'' büyüdüğünü düşünürüz. Ağacın en üstteki (yaprak)
düğümleri aksiyomlardır ve \emph{başlangıç sekantları} diye adlandırılır.
En alttaki sekanta \emph{son sekant} denir. $\pi$, son sekantı~$S$ olan bir
!!a{proof} ise $\pi \Proves S$ de yazarız. Bir $S$ sekantının bir
!!a{proof}{}ı varsa $\Proves S$ yazarız; istenirse ispatın bulunduğu
sistem $\Proves$ simgesinin solunda belirtilir, örneğin $\Log{G1c} \Proves
!A, !B \Sequent !A \land !B$. Başlangıç sekantları dışında bir !!a{proof}
içindeki her sekant, bir~$R$ kuralıyla yapılan çıkarımın \emph{sonucudur}.
Bir !!a{proof} içinde bir sekantın hemen üstünde duran bir ya da iki sekanta
(yani düğümün ebeveynlerine) çıkarımın \emph{öncülleri} denir.

Bir !!a{proof}{}ın tümevarımsal kuruluşunda kullanılan !!{proof}s, onun
\emph{alt-!!{proof}s{}ı} diye adlandırılır. Bir çıkarımın öncülleriyle biten
alt-!!{proof}s, çıkarımın \emph{doğrudan alt-!!{proof}s{}ı} diye adlandırılır.

!!{proof}s{}ın karmaşıklığını çoğu kez bir doğal sayıyla ölçmek gerekir.
Bir !!a{proof} içindeki sekantların sayısını doğrudan sayabilirdik; ancak
çoğu zaman daha yararlı ölçü, bir !!a{proof}{}ın !!{height} değeridir: son sekant
ile bir başlangıç sekantı arasındaki en büyük sekant sayısıdır. Bunun da
tümevarımsal tanımını verebiliriz.

\begin{defn}
Bir !!a{proof}~$\pi$'nin \emph{!!{height} değeri}~$\pheight{\pi}$ tümevarımsal
olarak şöyle tanımlanır:
\begin{enumerate}
  \item $\pi$ tek bir aksiyomdan oluşuyorsa $\pheight{\pi} = 0$.
  \item $\pi$, doğrudan alt-!!{proof}{}ı~$\pi_1$ olan tek öncüllü bir
  çıkarımla bitiyorsa $\pheight{\pi} = \pheight{\pi_1} + 1$.
  \item $\pi$, doğrudan alt-!!{proof}s{}ı~$\pi_1$ ve~$\pi_2$ olan iki
  öncüllü bir çıkarımla bitiyorsa $\pheight{\pi} =
  \max(\pheight{\pi_1}, \pheight{\pi_2}) + 1$.
\end{enumerate}
Bir $S$ sekantının yüksekliği $\le n$ olan bir !!a{proof}{}ı varsa
$\Proves[n] S$ yazarız.
\end{defn}

\begin{ex}
\Log{G3c}'de $!C$ atomik olmak üzere $!C, \Gamma \Sequent \Delta, !C$
biçimindeki her sekant bir aksiyomdur. $!D$ ile $!E$'nin atomik tümceler
olduğunu varsayalım. O zaman $!D, !E \Sequent !D$ ve $!D, !E \Sequent !E$,
$!C, \Gamma \Sequent \Delta, !C$ aksiyomunun örnekleridir. Örneğin $!C,
\Gamma \Sequent \Delta, !C$ içinde $!C \ident !E$, $\Gamma = \{!D\}$ ve
$\Delta = \emptyset$ alırsak tam olarak $!D, !E \Sequent !E$ sekantını
elde ederiz. (Çoklu\emph{kümeler}le çalıştığımızı; dolayısıyla $!D, !E
\Sequent !E$ ile $!E, !D \Sequent !D$'nin aynı sekant olduğunu anımsayın.)

$!D, !E \Sequent !D$, \Log{G3c}'nin bir aksiyomunun örneği olduğundan tek
başına \Log{G3c}'de bir !!a{proof} sayılır; $!D, !E \Sequent !E$ de
\olref{defn:proof}\olref{defn:prf:axiom} gereğince öyledir. Bunlara
$\pi_1$ ve~$\pi_2$ diyelim. \olref{defn:prf:two-prem} gereğince bu iki
!!{proof}s{}ı alıp iki öncüllü~$\RightR{\land}$ kuralını kullanarak
yeni bir ispat ($\pi_3$) üretebiliriz:
\[
\Axiom$!D, !E \fCenter !D$
\Axiom$!D, !E \fCenter !E$
\RightLabel{\RightR{\land}}
\BinaryInf$!D, !E  \fCenter !D \land !E $
\DisplayProof
\]
Ardından $\pi_3$'ten, tek öncüllü~$\LeftR{\land}$ kuralını kullanarak
(\olref{defn:proof}\olref{defn:prf:one-prem}) şu !!{proof}~$\pi_4$'ü elde
ederiz:
\[
\Axiom$!D, !E \fCenter !D$
\Axiom$!D, !E \fCenter !E$
\RightLabel{\RightR{\land}}
\insertBetweenHyps{\hspace{5ex}}
\BinaryInf$!D, !E  \fCenter !D \land !E $
\LeftSubproofLabel{$\pi_3$}
\RightLabel{\LeftR{\land}}
\UnaryInf$!D \land !E  \fCenter !D \land !E$
\RightSubproofLabel{$\pi_4$}
\DisplayProof
\]
$\pi_4$'teki son çıkarımın doğrudan alt-!!{proof}{}ı~$\pi_3$'tür.
$\RightR{\land}$ çıkarımının doğrudan alt-!!{proof}s{}ı $\pi_1$ ve~$\pi_2$'dir.
$\pi_4$'ün alt-!!{proof}s{}ı $\pi_1$, $\pi_2$, $\pi_3$ ve $\pi_4$'ün
kendisidir. $\pheight{\pi_1} = \pheight{\pi_2} = 0$,
$\pheight{\pi_3} = 1$ ve $\pheight{\pi_4} = 2$ olur. Her atomik $!D$ ve
$!E$ için $\Log{G3c} \Proves[4] !D \land !E \fCenter !D \land !E$
olduğunu göstermiş bulunuyoruz.
\end{ex}

